(1) 两无限大平行电极相距d，阴极电势0，阳极电势V>0；电子质量m、电荷-e，以可忽略初速度由阴极发射，一维定常无碰撞流且阴极场为零。由能量守恒、连续性和Poisson方程建立电势phi(x)的微分方程。
(2) 积分微分方程得到phi(x)、速度v(x)、电荷数密度n(x)及可传输电流密度绝对值j的Child–Langmuir律。
(3) 求粒子阴极到阳极的渡越时间，核验电荷驻留量/面积与j乘渡越时间相等。
(4) 比较外加电压或极板距的标度，说明阴极奇异密度、非零注入速度和有限温度导致的适用边界。
